\chapter{案例研究与专家评估}\label{chap:case}

本章把前述系统放回历史学家的工作场景，围绕一项案例研究、三个使用场景与一项小规模专家评估，验证"图模型—智能体—人工"三条证据流能否协同收敛。数字均来自原型部署的实测记录。

\section{两阶段案例：先验抽取与跨队列迁移}\label{sec:case-zhang}

案例由一位长期使用 CMGPD-LN、专攻清代辽东人口史的"张教授"完成，其关切是：某一队列中观察到的婚姻规范能否以可复用先验的形式跨队列迁移。

\textbf{校准阶段（1882 队列）。}时间滑块定位至 1882 年后，蜂窝画布显示已确认婚姻聚集成数个空间簇，提示嵌入空间已沿潜在社会分层组织。张教授选取中等密度格点 \texttt{\#136}，其中含两对真婚配 \texttt{P23497--P24248} 与 \texttt{P77574--P79028}。为 \texttt{P23497} 启动六轮辩论后，他注意到双方智能体反复以"共享八旗"与"父系延续"为核心论据，对官职暗示的家境讨论较少，这与清代旗内婚的制度性张力高于经济相称性的既有认知吻合。他随即把父系权重提升至~1.65、共享八旗权重提升至~1.55；对属表亲层级、得分边际较窄的 \texttt{P77574}，刻意保持权重于 1.10–1.30 的中等区间。经数轮重协商，两位真妻子均被推至第一名；同一组先验服务结构不同的两段婚姻，使他确信先验捕捉到真实局部规范。点击"保存为格点先验"后，聚合权重（父系~1.15、同胞~1.14、同户~1.12、八旗~1.25）写入格点 \texttt{\#136}。

\textbf{部署阶段（1885 队列）。}前移至 1885 年后，画布切换至部署态等高线，格点 \texttt{\#136} 填入 11 位全新丈夫候选，呈深蓝高置信，叠于 1882 年密度等高区之上，金色规则点持续可见。选取 \texttt{P91883} 启动协商，双方智能体在第一轮即自发把旗内婚与世系结构作为讨论中心，与 1882 年的论据走向一致；先验导引关注重心却未替智能体作结论。处理完格点全部候选后，整体协商较校准阶段收敛更快、人工干预更少。张教授由此判断：迁移的并非"谁与谁结婚"的个体记忆，而是"哪些证据值得最先被讨论"的群体级空间经验。他据此小结两点：八旗与父系主导的婚姻规范在 1882 至 1885 年间持续稳定，与清代制度史共识吻合；把先验空间锚定、再沿时间迁移具有可推广的方法论意义。

\section{三个使用场景}\label{sec:case-uc}

\textbf{UC1：批量提交后的智能体仲裁。}专家 E1 在 1882 队列上半象限发现一簇绿色高置信、属同户骨架的蜂窝格，将批量阈值 \(\theta\) 设为 1.0 后于 90 秒内提交 8 项推荐中的 7 项。对剩余不确定候选，他启动协商竞技场；智能体在事件登记中检索到丈夫一侧带有"逃离八旗登记"记录，E1 随即注入 \texttt{@everyone be skeptical of stability claims}，二者在第四轮均下调评分。E1 拒绝该候选，并撤回了一笔早先的批量提交，最终以 HGT 与 MAS 双来源标记结案。

\textbf{UC2：消融驱动的方法学审计。}专家 E3 在 1890–1895 队列上交替开关母题级消融。当 \(\mathsf{m}_4\)（旗内婚）被消融后，蜂窝格点结构显著变化：未消融下进入顶五分位的三对候选在去除八旗证据后跌至阈值线下。E3 把这一现象记录为旗内婚的"证据混淆"，并对受影响候选悉数搁置。

\textbf{UC3：智能体引导的姓氏推理。}专家 E2 在 1885 队列注意到一位王姓候选被异构图变换器排在 12 位中第六位且无显著母题匹配。他注入 \texttt{@everyone wife's surname is Wang}，唯有该候选的妻方智能体显式给出姓氏内婚论据，得分由 4.2 升至 7.8。E2 接受该候选，以"仅 MAS"标记提交并入审计待复核队列。

\section{专家评估}\label{sec:case-expert}

我们邀请 5 位领域专家（3 位历史学者、2 位计算家谱学者）于 1882 与 1903 两个队列上完成 75 分钟的结构化访谈式走查。访谈结束后专家在 1–7 分 Likert 量表上对三项指标作整体评分，结果汇总于表~\ref{tab:expert-likert}。

\begin{table}[h]
  \centering
  \caption{专家评估 Likert 评分（1--7 量表，$n=5$）。}
  \label{tab:expert-likert}
  \begin{tabular}{lc}
    \toprule
    指标维度 & 均值（标准差）\\
    \midrule
    可用性（usefulness） & 6.0~(0.71) \\
    可解释性（interpretability） & 6.0~(0.71) \\
    对仅 MAS 来源提交的信任 & 5.0~(1.00) \\
    \bottomrule
  \end{tabular}
\end{table}

开放问题归纳出三项定性发现。其一，提示控制台是最一致赞许的特性，5 位专家均将其视为"主观能动感提升最显著"的功能，2 位专家建议增设 \texttt{@household-head} 路由标签。其二，来源标记切实改变了阅读行为：在同时携带 HGT 与 MAS 双标记的 23 例提交中，19 例无须展开协商转录即被接受；而在仅 MAS 标记的 14 例中，11 例触发了对完整转录的展开阅读，说明专家把"仅 MAS"视为需追加证据投入的信号。其三，撤回功能被实质使用：所有会话中累计点击 14 次撤回，其中 9 次最终改变承诺——或接受不同候选，或选择留空——这印证设计目标 DG5 的可逆性是真实运转的功能而非安全网，也与 Sacha 等\citep{sacha2014knowledge} 提出的知识生成模型中"假设—证伪—修订"的节律相符。

最保守专家在"对仅 MAS 来源提交的信任"上仅给出 4.0 分，把均值拉至 5.0；这一保留态度与 Turpin 等\citep{turpin2023faithfulness} 关于大语言模型解释可能是事后合理化的警示一致：缺少图模型旁证时，单凭智能体推理作婚姻补全判断仍应审慎对待。综上，本系统的可视分析闭环既降低了批量推荐的复核成本，又在不确定个案上把责任清晰分摊给人、模型与智能体三方。
